\documentclass[12pt]{article}
\usepackage[a3paper, margin=1in]{geometry}
\usepackage{amsmath, amssymb, amsthm, ulem, booktabs}

\begin{document}

\section*{Domácí úkol 1.2}

Vypracovali: Daniel \textit{``Localhost"} Ransdorf, Jan \textit{``kokeshh"} Kodeš \hfill Podpis: \rule{4cm}{0.4pt}

\bigskip

Matice A je symetrická, označme tedy:

\[
  \left( \begin{array}{cc} a&b\\b&d \end{array} \right) := A
\]

Pro aritmetický vektor $\vec{g} \in \mathbb{R}^2$, označme jeho souřadnice: $\left(\begin{array}{c} g_1 \\ g_2 \end{array}\right) := \vec{g}$

\begin{enumerate}
  

























  \item $\Leftarrow$: Když $a>0,\ detA>0$, pak je $<,>$ skalární součin.
    \medskip

    Předpokládejeme $a>0,\ detA>0$, pak platí $ad-b^2 > 0 \ (*)$.
    \[
      ad-b^2 > 0 \quad \Rightarrow \quad ad>b^2\geq0 \quad \Rightarrow \quad d>0 \quad (**)
    \]

    Když dokážeme, že je matice $A$ pozitivně definitní, pak je $<\vec{x},\vec{y}> = \vec{x}^T A \vec{y}$
    nutně skalární součin podle tvrzení ze skript.

    Reálná je pozitivně definitní, pokud:
    \begin{enumerate}
      \item je hermitovská
      \item $\vec{x}^T A \vec{x} \geq 0$
      \item $\vec{x}^T A \vec{x} = 0 \quad \Longleftrightarrow \quad \vec{x} = \vec{0}$
    \end{enumerate}

    Dokažme podmínky:

    \begin{enumerate}
      \item Pracujeme na realnými čísly, čili tento bod plyne ze symetrie matice.
      \item Zvolme libovolný vektor $\vec{x} \in \mathbb{R}^2$, počítejme dle zadaného vzorce:
      
        \[
          \vec{x}^T A \vec{x} = \vec{x}^T \left( \begin{array}{cc}
            ax_1 & +bx_2 \\ bx_1 & +dx_2
          \end{array} \right) = ax_1^2 + bx_1x_2 + bx_2x_1 + dx_2^2 = ax_1^2 + 2bx_2x_1 + dx_2^2
        \]
        Zvolme libovolně $x_2\in\mathbb{R}^2$, zkoumejme parabolu podle $x_1$. Rozdělme případy:
        \begin{enumerate}
          \item $x_2 \neq 0$:
            \begin{align*}
              (a)x_1^2 + (2bx_2)x_1 + (dx_2^2) \\
              D = 4b^2x_2^2 - 4adx_2^2 = \underbrace{4x_2^2}_+\underbrace{(b^2-ad)}_{- \text{ dle }(*)} < 0
            \end{align*}
            Z diskriminantu plyne, že neexistuje $x_1$ takové, že součin bude roven nule. 
            A protože platí $a>0$, parabola se nachází celá nad 0. Tedy $\vec{x}^T A \vec{x} > 0$
          \item $x_2 = 0$:
            \[
              (a)x_1^2 + (2bx_2)x_1 + (dx_2^2) = (a)x_1^2 +0+0 = (a)x_1^2
            \]
            Jelikož $a>0$, tato parabola je nezáporná a je rovna nule právě tehdy, když $x_1=0$.
            Čili $\vec{x}^T A \vec{x} \geq 0$, přičemž $\vec{x}^T A \vec{x} = \vec{0} \ \Leftrightarrow \ x_1=0$.
        \end{enumerate}

        Každopádně dostáváme $\vec{x}^T A \vec{x} \geq 0$.
      
      \item Z úvah b.i, b.ii plyne:
        \[\vec{x}^T A \vec{x} = 0 \quad \Longleftrightarrow \quad \vec{x} = \vec{0}\]
    \end{enumerate}

    Tedy matice A je pozitivně definitní, a tedy $<,>$ je skalární součin. $\quad \quad \square$

  \pagebreak
  
  \item $\Rightarrow$: Když je $<,>$ skalární součin, pak platí $a>0,\ detA>0$.
    \medskip
    
    Předpokládejme, že $<,>$ je skalární součin. Dle tvrzení ze skript víme, že $A$ je pozitivně definitní.
    Víme tedy:

    \begin{itemize}
      \item $\vec{x}^T A \vec{x} \geq 0 \quad (*)$
      \item $\vec{x}^T A \vec{x} = 0 \ \Leftrightarrow \ \vec{x}=\vec{0} \quad (**)$
      \item $\vec{x}^T A \vec{x} > 0 \ \Leftrightarrow \ \vec{x} \neq \vec{0} \quad (***)$
    \end{itemize}

    Rozepišme součin jako v bodě 1:

    \[
      \vec{x}^T A \vec{x} = ax_1^2 + 2bx_2x_1 + dx_2^2 \overset{(***)}{\geq} 0
    \]

    Zvolme libovolně $x_2\in\mathbb{R}$, zkoumejme parabolu podle $x_1$. Rozdělme případy, čemu může být rovno $x_2$:

    \begin{enumerate}
      \item $x_2 = 0$:
        \[ax_1^2 + 2bx_2x_1 + dx_2^2 = ax_1^2\]
        Kdyby $a=0$, pak bychom zvolením $x_1=1$ získali spor s $(***)$: $\quad a\cdot1 = 0 \overset{(***)}{>} 0$
        
        Kdyby $a<0$, pak bychom zvolením $x_1=1$ získali spor s $(*)$: $\quad 0> a = a\cdot(1) \overset{(*)}{\geq} 0$.

        Tedy nutně platí \underline{\underline{$a>0$}}.
      \item $x_2 \neq 0$:
        Určeme diskriminant kvadratické rovnice. $\vec{x} \neq \vec{0}$, čili podle $(***)$ nemůže existovat $x_1$ takové, že
        bude výraz rovný nule.
        A tedy diskriminant musí být záporný.
        \begin{align*}
          &(a)x_1^2 + (2bx_2)x_1 + (dx_2^2) \\
          D &= 4b^2x_2^2 - 4adx_2^2 = \underbrace{4x_2^2}_+(b^2-ad) < 0 \\
          \Longrightarrow \quad &b^2-ad < 0 \quad \Longrightarrow \quad ad-b^2 > 0 \quad \Longrightarrow \quad detA>0
        \end{align*}

        Tedy nutně platí \underline{\underline{$detA>0$}}.
    \end{enumerate}

    Z předchozích bodů víme, že platí $a>0, detA>0$. $\quad\quad \square$
    

    
















































\end{enumerate}

\end{document}
